% =====================================================================
% DRAFT (10 Oct 2026) -- Section: evaluation protocol
% Every number carries a "% fonte:" comment. Missing numbers are \todo{}.
% Labels referenced from sec_gt.tex: sec:protocol_baselines,
% sec:trivial_baselines, sec:protocol_human.
% =====================================================================

\section{Evaluation Protocol}
\label{sec:protocol}

The protocol below was fixed before the corresponding numbers were computed.
Each experiment has a written protocol whose SHA-256 hash was recorded before
the first job ran; later changes are dated amendments, also hashed before any
new number exists.
% fonte: aau/runs/evidence/e12/protocol.md (+ protocol.sha256, protocol_amend1.sha256); aau/runs/evidence/trainer_v3/factorized_protocol.md (+ .sha256 ed emendamenti 1-3); aau/human_study_v2/PROTOCOL.md (+ PROTOCOL.sha256)
We release the protocols with their hashes \todo{decidere se rilasciarli nel materiale supplementare}.

\subsection{Domains and Roles}
\label{sec:protocol_domains}

Table~\ref{tab:domain_roles} lists every data source and its role. Three rules
apply. (i) Test domains are never used for training or for model selection;
the data pipeline refuses them by construction (loading HIFI3D, FaceVerse or
FaceScape as a training source raises an error).
% fonte: aau/runs/evidence/stream/massive_ready.md, sez. 1 ("Domini vietati: FaceScape, FaceScape50, HIFI3D e FaceVerse sollevano RoleError")
(ii) Hyperparameters are chosen only on the development domain, FaceScape,
with a score declared in advance (the mean of the graded Spearman correlation
without crop and the rank-1 accuracy with expressions); the reported
checkpoint is the last exponential-moving-average checkpoint, with no
selection.
% fonte: paper/PLAN_MASSIVE.md, sez. 9 (Dev, Punteggio dev, Checkpoint finale) e sez. 14.5 (FaceScape solo dev, mai in training)
\todo{dichiarare con quale GT si calcola il punteggio dev (FR?) per il run finale}
(iii) We distinguish \emph{untouched} test sets from test sets that we
inspected during development. HIFI3D and FaceVerse were evaluated repeatedly
while the method was developed (37 configurations on HIFI3D alone) and are
therefore reported as \emph{seen during development}; the untouched tests are
the FaMoS test subjects, NoW and Multiface.
% fonte: paper/PLAN_MASSIVE.md, sez. 14.5 ("HIFI3D, con 37 bracci gia' valutati, e FaceVerse sono in pratica 'visti durante lo sviluppo' e vanno dichiarati cosi'")
One caveat applies to FaMoS: all 95 FaMoS subjects, test subjects included,
entered the identifiability arbiter that selected the reference
(Sect.~\ref{sec:gt_arbiter}). No model was trained or selected there, but the
choice of FR is not independent of these 15 people.
% fonte: aau/runs/evidence/e12/protocol.md, Emendamento 1, E2 ("Persone: tutte le 95 registrate (TRAIN + TEST)")

\begin{table}[!ht]
\centering
\small
\caption{Data sources and roles. Training sources follow the configuration
of the final run \todo{confermare dopo la revisione con l'utente del run
massivo}. ``Seen'' test sets were inspected during development.}
\label{tab:domain_roles}
% fonte: aau/runs/evidence/stream/massive_ready.md, sez. 1 (preset massive: bfm2019, ict, gnm, flame2023, famos 80 TRAIN); paper/PLAN_MASSIVE.md, sez. 2, 9, 14.5
\begin{tabular}{llll}
\toprule
Source & Type & Units & Role \\
\midrule
BFM 2019~\cite{bfm2019}            & 3DMM          & mm (declared)       & training \\
ICT-FaceKit (Light)~\cite{ict-facekit} & 3DMM      & cm (declared)       & training \\
GNM Head~\cite{gnm}                & 3DMM          & m (declared)        & training \\
FLAME 2023 Open~\cite{flame}       & 3DMM          & m (declared)        & training \\
FaMoS, 80 subjects~\cite{famos}    & real, FLAME registrations & mm      & training \\
\midrule
FaceScape bilinear~\cite{facescape} & 3DMM from real scans & mm (declared) & development only \\
\midrule
HIFI3D~\cite{hifi3d}               & 3DMM          & IPD-calibrated      & test (seen) \\
FaceVerse~\cite{faceverse}         & 3DMM, expressions & IPD-calibrated  & test (seen) \\
FaMoS, 15 subjects~\cite{famos}    & real scans and registrations & mm   & test (untouched) \\
NoW validation~\cite{now}          & real images, scans & mm (scans)     & test (untouched) \\
Multiface~\cite{wuu2022multiface}  & real, tracked meshes & mm           & test (untouched) \\
\bottomrule
\end{tabular}
\end{table}
% fonte unita': aau/runs/evidence/e12/summary.md, sez. 1 (unita' dichiarate e IPD; HIFI3D 9.34 mm/unita', FaceVerse 161.8 mm/unita')
% fonte BFM 2019 in mm: aau/runs/evidence/stream/massive_ready.md, sez. 2 ("BFM 2019: mm dichiarati")

\paragraph{Splits and near-duplicates.}
The FaMoS test subjects are subjects 079--093, the 15 people with test scans,
fixed before use; the closest pair between FaMoS training and test subjects is
2.17\,mm apart in the common region.
% fonte: paper/BOARD_DIARY.md, "9 ottobre: FaMoS preparato" (split TEST = 079-093, distanza minima 2.17 mm); aau/famos/split.json
On HIFI3D no test identity has a training neighbour closer than the closest
pair inside the HIFI3D pool itself; GNM is, however, systematically the
training domain closest to HIFI3D (median nearest-neighbour distance 2.73\,mm,
against 2.79\,mm for ICT, 3.45\,mm for BFM and 2.54\,mm inside the HIFI3D pool).
% fonte: aau/runs/evidence/e8/summary.md, Conclusioni 5
The means of the domains form two groups 3.2--3.8\,mm apart: FLAME, BFM, ICT
and Multiface on one side, FaceScape, HIFI3D and FaceVerse on the other. The
second group consists of models built on East-Asian populations; this is a
hypothesis about the cause, which we did not verify.
% fonte: aau/runs/evidence/e8/summary.md, Conclusioni 6
Since FaceScape lies 0.84\,mm from HIFI3D, it is used for development only
and never for training.
% fonte: aau/runs/evidence/e8/summary.md, Conclusioni 6 ("FaceScape dista 0.84 mm da HIFI3D"); paper/PLAN_MASSIVE.md, sez. 14.5

\paragraph{Topologies.}
On the 3DMM test and development domains each identity is realized in the six
topologies of REMESH (Sect.~\ref{sec:remesh_benchmark}): original, noisy,
remesh, crop, down8k and up60k. Graded results use 100 identities per domain
(drawn with a fixed seed from a pool of 500) and all cross-topology pairs of
different subjects; the primary group is \emph{no-crop cross}, the 20 ordered
pairs of distinct topologies without crop. We also report all 30 ordered
pairs and the per-subject-pair mean, and crop separately.
% fonte: aau/runs/evidence/e12/protocol.md, sez. 4 (100 soggetti di select_subjects(...,1234); nocrop_cross 20 coppie ordinate PRIMARIO, all_cross 30, subject_pair_mean); aau/runs/competitors_hifi3d/summary.md, Dati (600 mesh = 100 x 6 topologie)
A seventh tessellation generator, never used in training, is held out as an
additional test topology \todo{stato del ``remesher mai visto'' (PLAN\_MASSIVE sez. 2.5): generatore, numeri}.
% fonte: paper/PLAN_MASSIVE.md, sez. 2 (Discretizzazioni, punto 5) e sez. 9

\subsection{Metrics}
\label{sec:protocol_metrics}

\paragraph{Graded distance.}
On domains with dense correspondence, the primary statistic is the Spearman
correlation between a candidate distance and the reference over cross-subject
mesh pairs (Eq.~\eqref{eq:spearman_correlation}). The reference is FR for
metric data and SR for scale-free data (Sect.~\ref{sec:gt_definition}); every
table also reports the legacy maxabs reference and the full Procrustes
reference so that the conclusions can be checked under each
\todo{colonna SR per tutti i metodi: non ancora calcolata su HIFI3D}.
% fonte: paper/PLAN_MASSIVE.md, sez. 19 ("(c) robustezza: le conclusioni del paper devono reggere con GT-F e con GT-S") e sez. 20 (decisione GT = FR)

\paragraph{Recognition.}
Recognition uses identity labels only and is therefore independent of any
reference distance. On HIFI3D we use the five topologies without crop: each
of the 2{,}000 queries is retrieved against a gallery of the 100 identities
in another topology (rank-1, and mAP, which here equals the mean reciprocal
rank), and verification uses 1{,}000 genuine and 99{,}000 impostor pairs
(AUC); crop is reported separately.
% fonte: aau/runs/data_scale_ood/arcface_vs_scale_hifi3d.md, sez. (b) ("5 topologie senza crop: 2000 query di retrieval (galleria di 100 mesh in un'altra topologia; mAP = MRR), verifica su 1000 coppie stessa persona e 99000 diverse")
We additionally report TAR at FAR $=10^{-3}$ and $10^{-4}$; with 99{,}000
impostors the latter rests on about ten impostors above threshold and is
noisy.
% fonte: aau/runs/indomain_recog/summary.md, Revisione 1, punto D-E (TAR@FAR 1e-3 e 1e-4; "la seconda e' rumorosa")
On FaceVerse the queries carry expressions and the gallery is neutral; on the
FaMoS test subjects, raw scans with expressions are retrieved against neutral
scans \todo{protocollo FaMoS test definitivo (scansioni grezze vs registrazioni)}.
% fonte: paper/PLAN_MASSIVE.md, sez. 9 (Metriche di test: rank-1, mAP, AUC, TAR@FAR); paper/BOARD_DIARY.md, "9 ottobre: FaMoS preparato" (riconoscimento scansione con espressione -> scansione neutra)
On NoW, where no dense correspondence exists, we compare the ranking of
reconstruction methods induced by each distance with the ranking induced by
the official NoW error \todo{definire la statistica (Kendall $\tau$) e il riferimento SR per NoW; aau/runs/now\_eval/protocol.md}.
% fonte: aau/runs/now_eval/summary.md (protocollo NoW: 3DDFA_V2, SynergyNet, PRNet; errore ufficiale)

\paragraph{Uncertainty and decision rules.}
All intervals are 95\% bootstrap intervals with 1{,}000 replicates that
resample \emph{subjects} (or persons, on FaMoS), not pairs; a pair is weighted
by the product of the multiplicities of its two subjects. Differences between
methods, or between references, are paired: both sides are computed on the
same replicates, and we report the interval of the difference together with
the fraction of replicates in which it is $\le 0$. Two methods whose paired
interval contains zero are reported as not distinguishable, not as
equivalent.
% fonte: aau/runs/evidence/e12/protocol.md, sez. 3-4 (bootstrap per soggetto, 1000 repliche, seme 1234, peso = prodotto dei conteggi; delta appaiati con P(delta <= 0)); paper/PLAN_MASSIVE.md, sez. 16.2 ("assenza di prova, non equivalenza")
Decision rules are declared with the protocol. For example, the question
whether a model trained on FR approaches pairwise NICP on HIFI3D is decided
by the fraction $g=(\rho-0.194)/(0.367-0.194)$ of the gap between the
size-blind encoder and NICP that it closes: it ``reaches NICP'' if
$\rho\ge0.289$ (the lower end of the NICP interval) and its interval contains
0.367 or lies above; it ``approaches NICP'' if $g\ge0.5$ and the lower end of
its interval exceeds 0.194; the verdict counts only if both training seeds
agree.
% fonte: aau/runs/evidence/trainer_v3/factorized_protocol.md, sez. "Regola"
Outcome: \todo{esito dei bracci factorized e ctrl-FR secondo questa regola}.

\subsection{Baselines and Nuisance Removal}
\label{sec:protocol_baselines}

We evaluate 22 distances on HIFI3D under every reference: pairwise rigid ICP
followed by Chamfer, pairwise rigid ICP followed by non-rigid ICP with
point-to-triangle distance (NICP P2Tri), non-rigid registration to a domain
template at enrolment followed by a per-vertex distance (NICP on template),
two Chamfer variants, ArcFace~\cite{deng2019arcface} on shaded renders and on
normal maps (three views), the spectral descriptors
ShapeDNA~\cite{reuter2006shapedna} (two truncations, two normalizations),
HKS~\cite{sun2009hks} and WKS~\cite{aubry2011wks}, the zero-shot 3D foundation
encoders Uni3D~\cite{uni3d} and OpenShape~\cite{openshape} (native frame and
the exact frame of their training as ablations), and four earlier learned
checkpoints.
% fonte: aau/runs/evidence/e12/summary.md, sez. 5 (15 righe in tabella + 7 solo nei csv: uni3d_rx90, openshape_rx90, openshape_yzswap, shapedna_k50_l1, shapedna_k100_l1, hks, wks); aau/runs/competitors_hifi3d/summary.md (definizioni dei competitori)
All hyperparameters of the competitors were fixed before any of their numbers
existed.
% fonte: aau/runs/competitors_hifi3d/summary.md, protocollo ("Nessun iperparametro si sceglie guardando le metriche")

\paragraph{Nuisance removal consistent with the reference.}
A fair comparison gives every method the nuisance removal that its reference
declares unobservable, and nothing more. Under FR, registration baselines
operate in millimetres with a rigid ICP without scale; under SR, they use a
similarity ICP. Any canonicalization applied to our model is also applied to
the geometric baselines. We do not canonicalize at test time, because a
per-mesh rigid canonicalization lowered the graded correlation on HIFI3D
($-0.074$) and was neutral elsewhere.
% fonte: paper/PLAN_MASSIVE.md, sez. 19 ("Baseline eque: ... Per GT-F: ICP rigido, senza scala, in mm. Per GT-S: ICP di similarita'"), sez. 14.6 (simmetria), sez. 15 (E2: Spearman -0.074 su HIFI3D)
The registration baselines currently reported (Table~\ref{tab:gt_reeval})
normalize each mesh by its maximum absolute coordinate before a rigid ICP,
and the template baseline ends with a Procrustes step that includes scale;
both therefore remove size and are consistent with SR rather than FR.
% fonte: aau/runs/indomain_recog/summary.md, Revisione 1 (faceBench: maxabs per mesh, prealline bbox, ICP rigido senza scala; punto B: Procrustes con scala); aau/runs/competitors_hifi3d/summary.md, Aggiunta 14:28
\todo{rieseguire Chamfer, ICP+Chamfer, NICP P2Tri e NICP su template in mm con ICP rigido senza scala (riga FR) e con ICP di similarita' (riga SR) su HIFI3D, FaceScape dev, FaceVerse, FaMoS test}

\paragraph{Pairwise versus enrolment baselines.}
Pairwise pipelines (ICP, NICP) estimate one transformation per pair and pay
the registration for every comparison; enrolment baselines (NICP on
template, ArcFace, our encoder) pay once per mesh and compare cheaply
afterwards (Sect.~\ref{sec:cost}). We keep the two families separate in every
table, since only the second scales to retrieval in large galleries.

\subsection{Trivial Baselines}
\label{sec:trivial_baselines}

To measure how much of a reference a single number explains, we report two
trivial distances: \emph{size only}, $|\log S_i-\log S_j|$ with $S_i$ the
centroid size of Sect.~\ref{sec:gt_definition}, and \emph{height only},
$|\log H_i-\log H_j|$ with $H_i$ the vertical extent of the common region in
millimetres.
% fonte: aau/runs/evidence/e12/protocol.md, Emendamento 1, E1 ("Baseline banali (obbligatorie)")
Both are oracles computed on the reference shapes, not on the observed
meshes.

\begin{table}[!ht]
\centering
\small
\caption{Spearman correlation of the oracle trivial distances with each
reference (100 identities; FaMoS: 15 test subjects). 95\% subject-level
bootstrap intervals.}
\label{tab:trivial}
% fonte: aau/runs/evidence/e12/summary.md, sez. 3 (righe "F + rigida robusta" (FaMoS: "F (forma, mm)", che include la rigida robusta), "maxabs (legacy)", "unificata"; colonne "solo taglia", "solo altezza")
\begin{tabular}{llccc}
\toprule
Domain & Trivial distance & FR & Full Procrustes & maxabs \\
\midrule
HIFI3D     & size only   & 0.736 [0.66, 0.80] & 0.094 [0.01, 0.18]  & 0.139 [0.04, 0.24] \\
           & height only & 0.591 [0.49, 0.69] & 0.212 [0.09, 0.35]  & 0.172 [0.07, 0.29] \\
FaceVerse  & size only   & 0.209 [0.11, 0.31] & 0.019 [-0.07, 0.11] & 0.024 [-0.06, 0.11] \\
           & height only & 0.100 [0.02, 0.19] & 0.092 [0.01, 0.18]  & 0.116 [0.01, 0.20] \\
FaceScape  & size only   & 0.464 [0.34, 0.57] & 0.116 [0.02, 0.22]  & 0.068 [-0.03, 0.17] \\
           & height only & 0.526 [0.41, 0.61] & 0.322 [0.21, 0.43]  & 0.186 [0.09, 0.29] \\
FaMoS test & size only   & 0.787 [0.55, 0.90] & 0.213 [-0.24, 0.57] & 0.058 [-0.29, 0.38] \\
           & height only & 0.543 [0.23, 0.72] & 0.134 [-0.21, 0.46] & 0.099 [-0.24, 0.40] \\
\bottomrule
\end{tabular}
\end{table}

Table~\ref{tab:trivial} shows that FR is strongly driven by size on the
domains where size varies (HIFI3D, FaMoS), and weakly where the model
compresses it (FaceVerse). This is expected from Eq.~\eqref{eq:size_shape}
and is the intended behaviour of a form reference, but it also sets a floor:
on HIFI3D an oracle size scalar exceeds every evaluated method under FR
(Table~\ref{tab:gt_reeval}). We therefore report, next to every method, the
same trivial distances estimated from the observed meshes (centroid size and
height of each realization, across topologies)
\todo{baseline banali misurate sulle mesh osservate, non oracolo: Spearman con FR e riconoscimento},
and on FaMoS their identifiability (size only: AUC 0.939; height only: 0.889;
Table~\ref{tab:gt_arbiter}).
% fonte: aau/runs/evidence/e12/summary.md, sez. 4, famos95_first (solo taglia 0.939, solo altezza 0.889)
A distance that does not exceed the mesh-based size-only baseline is not
measuring shape beyond size.

\subsection{Human Study}
\label{sec:protocol_human}

The perceptual arbiter is a two-alternative forced-choice study (``Which face
is more similar to the reference face?'') designed to test whether size
matters for perceived similarity, which a render pipeline that normalizes each
mesh cannot test. The earlier version of the study rendered every face after
per-mesh maxabs normalization, so participants saw faces of equal height; it
is biased towards maxabs and cannot arbitrate size.
% fonte: paper/PLAN_MASSIVE.md, sez. 19, Arbitri (b) ("i render attuali normalizzano ogni mesh con maxabs ... lo studio cosi' e' sbilanciato verso maxabs")
The new version uses 100 GNM Head identities rendered at absolute scale: the
same orthographic camera and magnification for all faces (2.327\,px/mm; the
rendered face height is linear in its height in millimetres, $r=0.99998$),
three views, uniform grey material and fixed light, with the same robust rigid
fit as FR and never a per-mesh scale.
% fonte: aau/human_study_v2/PROTOCOL.md, sez. 1 (GNM Head v3.0, 100 identita', 2.327 px/mm, r = 0.99998, viste 0/45/90 gradi, rigida robusta, mai una scala per mesh)
Triplets are drawn from three strata in which two references disagree, each
with relative margin $\ge0.10$: FR vs.\ S (primary, 200 triplets), FR vs.\
size at nearly equal size (80), and S vs.\ maxabs at nearly equal size (80).
Each session has 2 practice trials, 18 test trials (12/3/3) and 2 attention
controls; a participant who fails a control is excluded.
% fonte: aau/human_study_v2/PROTOCOL.md, sez. 3-5
The statistic in each stratum is the share $q$ of responses agreeing with the
first reference, tested against $q=0.5$ with a crossed participants$\times$triplets
variance estimate~\cite{owen2007pigeonhole}; strata are tested in a fixed
order (gatekeeping), each at a nominal level of 0.04 calibrated by simulation
to keep the empirical type-I error at or below 0.05 in all 84 simulated
conditions.
% fonte: aau/human_study_v2/PROTOCOL.md, sez. 5 (Owen 2007, test a cascata, soglia tarata 0.04, 84 condizioni)
With the target of 60 retained participants, the primary stratum has 80\%
power for $q\ge0.60$ under the conservative heterogeneity setting; smaller
effects are reported with their interval only.
% fonte: aau/human_study_v2/PROTOCOL.md, sez. 5-6 (obiettivo 60 partecipanti; 80% per differenziali >= 0.20 su F_vs_S)
The primary stratum cannot distinguish FR from size alone, which agrees with
FR in all its triplets; FR is supported only if both the first and the second
stratum exceed 0.5.
% fonte: aau/human_study_v2/PROTOCOL.md, sez. 0 (avvertenze sulla lettura)
Declared limitations: one synthetic domain; the size contrast in the second
stratum (2--5\%) may be below visual threshold on small screens.
% fonte: aau/human_study_v2/PROTOCOL.md, sez. 3 (visibilita' del contrasto) e sez. 7
Participants recruited and excluded: \todo{N tenuti / N esclusi}; results:
\todo{q con IC per i tre strati}.
